%\resizebox{\textwidth}{!}{%
\begin{tikzpicture}

% Title
\node[dia/header=14cm] (title) {ПРОБЛЕМНО-ТЕМАТИЧНІ КЛАСТЕРИ};

% 2x2 grid of cluster boxes
\node[dia/rbox=5.5cm, below=1.5cm of title.south, xshift=-4cm] (c1)
  {Загальнотеоретичний і дидактичний кластер};
\node[dia/rbox=5.5cm, below=1.5cm of title.south, xshift=4cm] (c2)
  {Освітні практики};
\node[dia/rbox=5.5cm, below=1.2cm of c1] (c3)
  {Галузеві, порівняльні та узагальнювальні підходи};
\node[dia/rbox=5.5cm, below=1.2cm of c2] (c4)
  {Професійна педагогічна підготовка};

% Tree structure: title → row 1, row 1 → row 2
\draw[dia/line] (title.south) -- (c1.north);
\draw[dia/line] (title.south) -- (c2.north);
\draw[dia/line] (c1.south) -- (c3.north);
\draw[dia/line] (c2.south) -- (c4.north);

% Convergence point
\coordinate (conv) at ($(c3.south)!0.5!(c4.south) + (0,-1.5cm)$);
\node[dia/label] at (conv) (convlabel) {Точка перетину};

% Lines from row 2 boxes to convergence (direct, no overlap)
\draw[dia/line] (c3.south) -- (convlabel.north);
\draw[dia/line] (c4.south) -- (convlabel.north);

% Lines from row 1 boxes to convergence (routed outside row 2)
\draw[dia/line] (c1.south west) -- ([xshift=-0.4cm]c3.south west) -- (convlabel.north west);
\draw[dia/line] (c2.south east) -- ([xshift=0.4cm]c4.south east) -- (convlabel.north east);

% Bottom box
\node[dia/rbox=6cm, below=1.2cm of convlabel] (inclusive)
  {Інклюзивний освітній вимір};

% Arrow from convergence to bottom box
\draw[dia/arrow] (convlabel.south) -- (inclusive.north);

\end{tikzpicture}%
%}
